\documentclass[twocolumn]{aastex631}

\usepackage{amsmath}
\usepackage{multirow}
\usepackage{natbib}
\usepackage{graphicx} 
\usepackage{aas_macros}

\begin{document}

Higher-derivative quantum corrections are essential components of scalar-tensor effective field theories (EFTs), yet they typically reintroduce the Ostrogradsky ghost instability that the classical theory was designed to evade. This paper resolves this fundamental tension by establishing a rigorous equivalence between two distinct criteria for theoretical consistency. We analyze a general DHOST theory augmented by Gauss-Bonnet and Weyl-squared operators with coefficients that are arbitrary functions of the scalar field and its kinetic term. We then pursue two independent paths: first, we derive a set of differential equations for these coefficients by demanding that the full action remain invariant under the protective gauge symmetry of the classical theory. Second, we perform a first-principles Hamiltonian analysis using the ADM formalism, deriving a separate set of conditions by imposing the primary and secondary constraints required to eliminate the ghost. Our central result is the proof that these two sets of conditions, one algebraic and one dynamical, are mathematically identical. This equivalence demonstrates that the gauge symmetry is the fundamental origin of Hamiltonian stability in the quantum-corrected theory and establishes the symmetry principle as a powerful and practical tool for constructing consistent, ghost-free gravitational EFTs without resorting to a full Hamiltonian analysis.

\end{document}
                